This project focuses on the design, development, and implementation of an
interactive engineering tool that integrates electrical and systems control
systems. The objective is to develop a functional and reliable system
that allows users to interact with various electrical components and seamlessly import
or create one-line diagrams within a specialized component GUI. The software will validate 
system topology, map zones of protection, and perform critical momentary DC battery load 
calculations. By identifying worst-case scenarios, the tool provides GFT engineers with 
actionable insights to ensure maximum substation reliability and safety.


The project provides an opportunity to apply theoretical concepts learned
through coursework to the creation of a practical engineering application. Key areas of
focus include system design, component integration, programming, testing,
and troubleshooting. The completed system is intended to demonstrate the
application of engineering methods to solve a real-world problem.


The GUI allows users to construct protection system schematics using a Drag-and-Drop Approach, identifying and visualizing zones of protection, performing momentary current calculations, and accessing a library of electrical components. In addition, the tool provides the ability to import an existing PDF diagram and use it as a background reference while recreating or analyzing the system. These features are intended to make the analysis process more clean, organized, interactive, and easier to understand.


The tool is divided into several primary sections: the Protection Zone Outliner, Momentary Calculation Tool, Component Library, Sandbox, and PDF Importer/Exporter function.
\subsection{Project Objectives}